
\documentclass{article}
\usepackage{colortbl}
\usepackage{makecell}
\usepackage{multirow}
\usepackage{supertabular}

\begin{document}

\newcounter{utterance}

\centering \large Interaction Transcript for game `codenames', experiment `risk\_low', episode 2 with reflexion{-}Qwen3.5{-}27B.
\vspace{24pt}

{ \footnotesize  \setcounter{utterance}{1}
\setlength{\tabcolsep}{0pt}
\begin{supertabular}{c@{$\;$}|p{.15\linewidth}@{}p{.15\linewidth}p{.15\linewidth}p{.15\linewidth}p{.15\linewidth}p{.15\linewidth}}
    \# & $\;$A & \multicolumn{4}{c}{Game Master} & $\;\:$B\\
    \hline

    \theutterance \stepcounter{utterance}  
    & & & \multicolumn{2}{p{0.3\linewidth}}{
        \cellcolor[rgb]{0.9,0.9,0.9}{
            \makecell[{{p{\linewidth}}}]{
                \texttt{\tiny{[GM$|$GM]}}
                {'hidden': {'team': ['fork', 'shoe'], 'innocent': ['row', 'scientist', 'match', 'cotton', 'king', 'buffalo'], 'opponent': [], 'assassin': []}, 'revealed': {'team': {'team': [], 'innocent': [], 'opponent': [], 'assassin': []}, 'opponent': {'team': [], 'innocent': [], 'opponent': [], 'assassin': []}}}
            }
        }
    }
    & & \\ \\

    \theutterance \stepcounter{utterance}  
    & & \multicolumn{4}{p{0.6\linewidth}}{
        \cellcolor[rgb]{0.9,0.9,0.9}{
            \makecell[{{p{\linewidth}}}]{
                \texttt{\tiny{[P1$\langle$GM]}}
                \texttt{You have attempted this same task before and did not succeed. These are the plans you wrote after those attempts.} \\
\texttt{Trial \#1: New plan: The previous attempt was aborted because the clue "harry" is a proper noun (referring to Harry Potter) and likely violates the rule that the clue cannot be a specific name or part of a word, or it was deemed too ambiguous/invalid by the game master's strict parser. Additionally, while the format was correct, the choice of a proper noun as a single{-}word clue is risky. Next time, I will strictly use common nouns or verbs that are semantically related to the target words without being proper nouns or containing parts of the target words. For the remaining words (fork, cat, pound, witch, shoe, india, genius), I will choose a safer, more direct common noun like "magic" (for witch, genius, maybe cat) or "animal" (for cat, maybe others) and ensure the format remains exactly "CLUE: <WORD>\textbackslash{}nTARGETS: <LIST>" with no extra text. Specifically, I will avoid proper nouns entirely and select a generic concept that clearly links to at least two remaining targets without triggering the opponent's words.} \\
\\ 
\texttt{Trial \#2: New plan: My previous strategy failed because I prioritized guessing single words to avoid risk, which allowed the opponent to catch up and eventually win. Specifically, the clue "wear" was too broad, causing my partner to guess the distractor "dress" instead of "shoe," wasting a turn. In the future, I will prioritize grouping remaining words into multi{-}word clues to maximize efficiency per turn. For the final pair "fork" and "shoe," I should have used a more specific clue like "cutlery" (if shoe wasn't included) or accepted the risk of a broader clue like "dining" to target "fork" while explicitly listing only "fork" in the targets to prevent the partner from guessing "shoe" if the association was weak. More importantly, I must avoid clues that strongly associate with distractors (like "wear" {-}> "dress") by choosing more precise semantic links or splitting the targets if a single strong clue is too ambiguous. I will aim to clear 2+ words per turn whenever possible to outpace the opponent.} \\
\\ 
\texttt{Trial \#3: New plan: My previous attempt failed because I only targeted one word ("fork") with the clue "eat," leaving "shoe" unguessed and allowing the opponent to win. The strategy of playing it safe by targeting a single word was too slow. Next time, I will prioritize finding a single clue that semantically links to *all* remaining team words to clear them in one turn. For "fork" and "shoe," I should have used a broader concept like "dining" (if shoe fits a specific context) or, more likely, accepted that they might not share a strong single link and instead used a clue that strongly evokes "shoe" (like "foot") while explicitly listing only "shoe" in the targets to avoid the "row" distractor, then immediately following up with a second clue for "fork" in the next turn. Crucially, I must ensure the TARGETS list includes every word I intend the partner to guess in that specific turn, and I will avoid clues that strongly associate with distractors like "row" (e.g., avoiding "boat" or "oar") or "match" (e.g., avoiding "fire" or "game"). I will aim to clear the board as fast as possible, even if it requires two quick turns, rather than stalling with single{-}word clues that leave the board open.} \\
\\ 
\\ 
\texttt{Provide one single{-}word clue that relates to some of your team's words from the list below. You can choose to how many and to which words your clue relates to.} \\
\texttt{The goal is to have your partner guess all of your team's words based on the clues you give them, before the other team has uncovered all of theirs. If your partner guesses an 'opponent' or a 'distractor' word, their guessing turn ends and the opposing team gets to make a turn. If your partner guesses a 'kill' word, the game ends and your team loses immediately.} \\
\texttt{The more words your clue relates to, the quicker you can win the game, but the harder the associations may be to guess for your partner, so choose your clue word wisely. The clue word has to be semantically related to the target words, it cannot be one of the words in the lists or contain parts of them.} \\
\\ 
\texttt{Always give your single{-}word clue and your comma{-}separated list of related target words in the following format and make your answers as short as possible, never include any other text than is required in this form:} \\
\\ 
\texttt{CLUE: <WORD>} \\
\texttt{TARGETS: <TARGETS>} \\
\\ 
\texttt{Your team words are: fork, shoe.} \\
\texttt{Your opponent's team words are: .} \\
\texttt{Distractor words are: row, scientist, match, cotton, king, buffalo.} \\
\texttt{The kill words are: .} \\
            }
        }
    }
    & \\ \\

    \theutterance \stepcounter{utterance}  
    & \multicolumn{4}{p{0.6\linewidth}}{
        \cellcolor[rgb]{0.9,0.9,0.9}{
            \makecell[{{p{\linewidth}}}]{
                \texttt{\tiny{[P1$\rangle$GM]}}
                \texttt{CLUE: foot} \\
\texttt{TARGETS: shoe} \\
            }
        }
    }
    & & \\ \\

    \theutterance \stepcounter{utterance}  
    & & & \multicolumn{2}{p{0.3\linewidth}}{
        \cellcolor[rgb]{0.9,0.9,0.9}{
            \makecell[{{p{\linewidth}}}]{
                \texttt{\tiny{[GM$|$GM]}}
                \texttt{foot} \\
            }
        }
    }
    & & \\ \\

    \theutterance \stepcounter{utterance}  
    & & & \multicolumn{2}{p{0.3\linewidth}}{
        \cellcolor[rgb]{0.9,0.9,0.9}{
            \makecell[{{p{\linewidth}}}]{
                \texttt{\tiny{[GM$|$GM]}}
                ['shoe']
            }
        }
    }
    & & \\ \\

    \theutterance \stepcounter{utterance}  
    & & & \multicolumn{2}{p{0.3\linewidth}}{
        \cellcolor[rgb]{0.9,0.9,0.9}{
            \makecell[{{p{\linewidth}}}]{
                \texttt{\tiny{[GM$|$GM]}}
                {'word': 'shoe', 'assignment': 'team'}
            }
        }
    }
    & & \\ \\

    \theutterance \stepcounter{utterance}  
    & & & \multicolumn{2}{p{0.3\linewidth}}{
        \cellcolor[rgb]{0.9,0.9,0.9}{
            \makecell[{{p{\linewidth}}}]{
                \texttt{\tiny{[GM$|$GM]}}
                \texttt{opponent has won} \\
            }
        }
    }
    & & \\ \\

    \theutterance \stepcounter{utterance}  
    & \multicolumn{4}{p{0.6\linewidth}}{
        \cellcolor[rgb]{0.95,0.95,0.95}{
            \makecell[{{p{\linewidth}}}]{
                \texttt{\tiny{[P2$\rightarrow$P1]}}
                \texttt{Post{-}episode Reflexion (not sent to GM):} \\
\texttt{New plan: My previous attempt failed because I only targeted one word ("shoe") with the clue "foot," leaving "fork" unguessed and allowing the opponent to win. The strategy of playing it safe by targeting a single word was too slow. Next time, I will prioritize finding a single clue that semantically links to *all* remaining team words to clear them in one turn. For "fork" and "shoe," I should have used a broader concept like "dining" (if shoe fits a specific context) or, more likely, accepted that they might not share a strong single link and instead used a clue that strongly evokes "shoe" (like "foot") while explicitly listing only "shoe" in the targets to avoid the "row" distractor, then immediately following up with a second clue for "fork" in the next turn. Crucially, I must ensure the TARGETS list includes every word I intend the partner to guess in that specific turn, and I will avoid clues that strongly associate with distractors like "row" (e.g., avoiding "boat" or "oar") or "match" (e.g., avoiding "fire" or "game"). I will aim to clear the board as fast as possible, even if it requires two quick turns, rather than stalling with single{-}word clues that leave the board open.} \\
            }
        }
    }
    & & \\ \\

\end{supertabular}
}

\end{document}
